\begin{textsample}{NEG Dim 1 – vem\_ed\_27 – Score 82.00 – vem\_ed\_27/t143.txt}  \label{ex:f1_neg_173}
JORNADA DE PCIs: Breve histórico Desde 2021, a Jornada de Projetos Colaborativos Internacionais (PCIs / Cesu) tem o objetivo de compartilhar em língua portuguesa, com a comunidade acadêmica, boas práticas em Intercâmbios Virtuais do tipo COIL, bem como iniciativas de Internacionalização em Casa realizadas no CPS.
A primeira edição da Jornada ocorreu em novembro de 2021.
Em 2022, foi realizada como parte das atividades do International Virtual Exchange Symposium – IVES Cesu 2022.
No evento on-line, gestoras pedagógicas regionais compartilharam suas experiências sobre a divulgação de PCIs / Cesu em suas unidades e citaram exemplos realizados nas Fatecs da Regional São Paulo Sul e Baixada Santista; nas unidades de Matão, São José do Rio Preto e Itu.
Houve ainda um depoimento sobre o PCI Administrativo realizado entre gestores do CPS, Inacap (Chile), Uniminuto (Colômbia) e Universidade de Aveiro (Portugal).
Em 2021, com base na experiência dos PCIs / Cesu, ARInter e Cetec criaram o Programa de Aprendizagem Colaborativa Internacional (ProCin), voltado ao ensino médio, com estudantes da Prepa UDEM (México) e de Etecs do CPS.
A implantação do ProCin contou com a capacitação oferecida pela equipe dos PCIs / Cesu, que certificou28 professores e 636 estudantes, entre brasileiros e mexicanos, que realizaram 14 projetos em 2021.
O número 15 de VEm detalha esse evento: https://publicacoescesu.cps.sp.gov.br/VEm/article/view/274/214 Em 5 de outubro de 2023, ocorreu a terceira edição, cuja cobertura completa foi publicada em VEm 20: https://publicacoescesu.cps.sp.gov.br/VEm/article/view/283/219 continuação O webinário foi realizado por professores de Fatecs e de outras quatro instituições de ensino superior brasileiras: PUC- Campinas (São Paulo), UFF (Rio de Janeiro), Unesp (São Paulo) e Unichristus (Ceará).
O evento foi organizado em três painéis temáticos: a longevidade dos Intercâmbios Virtuais; a extensão desses projetos com a comunidade e o ensino de língua estrangeira no contexto dos Intercâmbios Virtuais.
ABERTURA Osvaldo Succi Junior explicou as diferentes nomenclaturas para os Intercâmbios Virtuais e destacou a “tropicalização” proposta para essa abordagem nos PCIs / Cesu do Centro Paula Souza.
Conhecidos em outros países como Collaborative Online International Learning (COIL), Global Learning Experience (GLE) ou Global Shared Learning (GSL), no Brasil também são chamados de BRaVE (Brazilian Virtual Exchange) pela Faubai, Associação Brasileira de Educação Internacional.
O Centro Paula Souza conquistou o selo BRaVE em outubro de 2022, fato noticiado em VEm 16: https://cesu.cps.sp.gov.br/wp-content/uploads/2023/02/Vem-16.pdf Succi Junior ressaltou a honra de receber contribuições de Instituições de Ensino Superior (IES) brasileiras para o evento, além dos trabalhos de professores de Fatecs, Universidade de Passo Fundo, Universidade Estadual de Maringá, Universidade Federal Rural de Pernambuco, Universidade Federal Fluminense, Universidade Vale do Rio Doce e Pontifícia Universidade Católica de Campinas.
Elenir Almeida Silva, representando a Coordenação de Línguas e Programas Internacionais da Cesu / CPS, deu as boas-vindas aos participantes. “É um prazer prestigiar este evento, que vem dar mais força ao sucesso dessa conexão entre Línguas e Internacionalização.
Esses projetos conectam pessoas, culturas de diversas partes do mundo e contribuem para a aprendizagem global que é o foco do Eixo de Línguas e Programas Internacionais da Cesu”.
Após a abertura, os participantes foram convidados a migrar para uma das três salas virtuais simultâneas disponíveis: Salas 1 e 2, “Boas Práticas”, e Sala 3, “Internacionalização em Casa: Ideias e Pesquisas”.
Um resumo das apresentações em cada uma das salas será apresentado nas páginas a seguir.
SALA 1 – BOAS PRÁTICAS Com moderação de Ana Carolina Freschi, responsável pelos projetos em inglês na equipe dos PCIs / Cesu e membro do conselho da Red LatAM COIL, a sala 1 trouxe apresentações sobre dois projetos realizados na Fatec Bragança Paulista – um deles, em três línguas, sob a liderança da professora de Espanhol Lilian de Souza. “Internacionalização Virtual: impacto dos Projetos Colaborativos Internacionais (PCIs) na formação e colocação profissional dos estudantes da Fatec de Bragança Paulista” foi o trabalho desenvolvido na Fatec Bragança Paulista pelos professores do Núcleo de Estudos da Linguagem da Fatec (Nelf) – Natalie Nara Mastrangi Goes, Ana Lúcia Leme Prestes, Lilian de Souza, Vera Encarnação Jordan de Aguiar – em colaboração com Paulo Henrique Leme Ramalho, professor de Informática na unidade. “Contagem Regressiva: 5,3,2,1: um PCI em 3 línguas”, de Lilian de Souza, sintetiza a experiência de envolver cinco professores trabalhando em três idiomas, com Instituições de Ensino Superior em dois países (Brasil e Chile) em um único projeto.
As equipes mistas integraram alunos dos cursos de Análise e Desenvolvimento de Sistemas (Brasil) e Introducción a Cloud Computing (Chile), que discutiram problemáticas presentes em ambos os países, tais como: saúde mental, manejo de resíduos, problemas com água, desmatamento e cuidados com pessoas idosas.
Para atividade final, foi proposto um site trilíngue (português, espanhol, inglês) para auxiliar a população em geral em relação às temáticas apresentadas. continuação Jancileidi Hübner, Cleonice Pletsch e Luciane Sturm (Universidade Passo Fundo) Jancileidi Hübner, Cleonice Pletsch e Luciane Sturm (Universidade Passo Fundo) destacaram o projeto COIL realizado na instituição como forma de estimular a formação docente.
Intitulada “Projeto COIL como Catalisador de Pautas Globais Urgentes na Formação Inicial de Professores”, a apresentação relatou a colaboração realizada entre a universidade gaúcha e a University of Akron (EUA).
O objetivo principal foi preparar os futuros professores para a integração de pautas transversais socioambientais nas práticas pedagógicas, por meio de tarefas síncronas e assíncronas que compuseram parte das avaliações de disciplinas curriculares.
Tais atividades “promoveram o debate intercultural sobre os Objetivos do Desenvolvimento Sustentável e a inserção dessas pautas no planejamento de aulas”, conforme descreveram as autoras no resumo do trabalho. continuação A importância dos Intercâmbios Virtuais na formação de professores também deu o tom à comunicação oral de Bruna Sampaio Silgueiro Mardegan e Luciana Cabrini Simões Calvo (Universidade Estadual de Maringá)."Reflexões Sobre Intercâmbio Virtual Na Formação Docente"abordou projetos realizados desde 2022 entre a UEM e Penn State University (EUA) e desde 2023 com a Universidad Pedagógica Nacional Francisco Morazán (Honduras).
No ano de 2024, o tema debatido foi a Inteligência Artificial e a Educação, em diálogo com várias Instituições de Ensino Superior (IES): além da universidade hondurenha já mencionada, as IES University of Erlangen-Nürnberg (Alemanha), Levinsky Wingate Academic College (Israel) e Universidade Federal do Espírito Santo participaram das discussões conjuntas.
SALA 2 – BOAS PRÁTICAS Pedagogia freireana na condução de projetos COIL entre Brasil e Equador; experiências e impressões sobre uma colaboração entre Fatec Indaiatuba e University of Minnesota Crookston (EUA); a realização de"Escrituras Narrativas"audiovisuais em projeto desenvolvido entre Fatec Ipiranga e Uniminuto (Colômbia) e um relato de COIL entre a Faculdade de Enfermagem da PUC-Campinas e universidades nos EUA e Filipinas.
Esses foram os temas abordados na"Sala 2 – Boas Práticas", que contou com moderação de Patrícia Patrício, responsável pela Comunicação do Eixo de Línguas e Programas Internacionais da Cesu / CPS.
Geyza Leyde Camello Lustosa, da UFRPE, começou a sessão trazendo a experiência realizada pela universidade pernambucana com a Universidad Nacional de Educação (Unae), no Equador.
A colaboração se deu por meio da Cátedra Paulo Freire, que visa estudar questões relativas à educação libertadora e à inclusão, juntamente com a Cátedra Aberta Latinoamericana e Caribenha da Unae. continuação Após formalizar acordo de cooperação, as IES fizeram um planejamento estratégico para desenvolver ações de Internacionalização em Casa ao longo de cinco anos (2023-2028).
Uma das ações foi o projeto COIL dentro da disciplina “Pedagogia de Paulo Freire”.
Os problemas principais a serem superados foram a grade curricular (dificuldade para reconhecer as atividades desenvolvidas nos projetos COIL dentro das disciplinas), a falta de regulamento para disciplinas on-line, a barreira linguística (português-espanhol) e as plataformas ou tecnologias acessíveis.
Por outro lado, as motivações envolvem fatores como fortalecimento de redes de cooperação, baixo custo para implementação das ações de Internacionalização em Casa, flexibilidade do COIL, interculturalidade, inclusão e envolvimento de disciplinas de língua estrangeira.
Os pontos positivos apontados na avaliação ao fim do projeto foram integração dos alunos, didática das professoras e material bilíngue.
Os aspectos a serem melhorados envolvem gravação das aulas síncronas (isso não foi feito) e superação das dificuldades com o idioma espanhol.
A professora de Inglês Talita Annunciato Rodrigues resumiu o projeto que orientou no curso de Comércio Exterior da Fatec Indaiatuba em colaboração com estudantes de Comportamento do Consumidor da University of Minnesota Crookston (EUA).
O tema tratado no projeto foi “Marketing Intercultural: considerações na abordagem de um mercado estrangeiro”.
Odenildo França Almeida, professor de Espanhol da Fatec Ipiranga, trouxe os principais resultados do PCI “Escrituras Narrativas” desenvolvido desde 2023 com a Uniminuto (Colômbia).
Estudantes de Tecnologia em Realização Audiovisual da universidade colombiana preparam curtas sobre temáticas urbanas contemporâneas, comuns à realidade de metrópoles como São Paulo e Bogotá, enquanto os estudantes da Fatec Ipiranga fazem a recepção crítica e dialogam sobre as produções dos colegas colombianos. continuação Yara Maria Randi apresentou o projeto COIL desenvolvido com as colegas Adeline Mariano Silva de Resende, Gabriela Marchiori Carmo Azzolin e Silvana Chorratt Cavalheri (PUC-Campinas), em colaboração com IES nos EUA e Filipinas.
O eixo norteador foi o desenvolvimento das competências linguísticas no idioma inglês dos futuros enfermeiros, já que essa é uma obrigatoriedade das diretrizes curriculares nacionais (2001).
Com a CSU-Chico (Califórnia, EUA), foram realizadas três edições do COIL entre 2023 e 2024, envolvendo 120 alunos brasileiros e 80 americanos na partilha de conhecimentos transculturais na área de enfermagem.
Com o San Pedro College (Filipinas), os modelos de liderança foram discutidos no segundo semestre de 2024 por 44 estudantes brasileiros e 66 filipinos.
Uma outra iniciativa de Internacionalização em Casa, com o condado de Sarasota (Florida), trabalhou traduções de carteiras de vacinação para facilitar assistência aos brasileiros nos EUA.
Além disso, a colaboração inclui pesquisas para posterior publicação acadêmica.
SALA 3 - INTERNACIONALIZAÇÃO EM CASA: IDEIAS E PESQUISAS Processos e práticas de institucionalização das ações de Internacionalização em Casa, bem como pesquisas derivadas de Intercâmbios Virtuais, deram a tônica à Sala 3, mediada por Osvaldo Succi Junior.
Vitor Ierusalimschy, chefe de Projetos Educacionais da Superintendência de Relações Internacionais da Universidade Federal Fluminense (UFF), ao comentar sobre “O COIL no Âmbito de uma Política Linguística para Internacionalização”, elencou as barreiras para a Internacionalização: Linguística - esta carrega em seu bojo a questão socioeconômica (dificuldade do acesso às línguas estrangeiras para as camadas mais pobres) e o desafio da inclusão das pessoas em vulnerabilidade por meio do ensino de idiomas estrangeiros nos cursos de Graduação.
Financeira - os custos para a mobilidade física são proibitivos para a maioria dos discentes das universidades públicas.
Nesse sentido, as ações de Internacionalização em Casa ajudam a dissolver essas barreiras.
Ierusalimschy ressaltou que os projetos COIL não vêm substituir a mobilidade tradicional, mas sim complementar.
Após essa explanação inicial, ele citou três casos de sucesso em sua instituição que aliam iniciativas como o Programa de Universalização de Línguas Estrangeiras (PULE) e projetos COIL.
COIL"PULE \& FIT"- realizado em 2021, em inglês, entre UFF e Fashion Institute of Technology (FIT, EUA).
Integrou formandos da universidade fluminense e ingressantes do instituto norte-americano, discutindo temas de identidade, cultura, pertencimento e ativismo político.
COIL"PULE \& ELAC"- realizado em francês com East Los Angeles College (ELAC, EUA), tinha o idioma como segunda língua para todos os envolvidos (brasileiros e norte-americanos).
Trocas culturais e diversidades regionais pontuaram as discussões.
PULE \& CSU Stan - realizado em português com estudantes da California State University Stanislaus (CSU Stan, EUA).
Para os brasileiros, as atividades desenvolvidas em dois encontros síncronos e três semanas de projeto trataram do ensino da língua portuguesa. continuação “Brazil and Armenia: a fruitful experience” trata do projeto desenvolvido entre as professoras Edilene Gasparini Fernandes (Fatec Rio Preto) e Anoush Ayunts (Yerevan State University, Armênia).
Na página 5 de VEm 20, é possível saber mais sobre a iniciativa: https://publicacoescesu.cps.sp.gov.br/VEm/article/view/283/219 Pedro Marçal e colegas da Universidade do Vale do Rio Doce (Univale) desenvolveram projeto COIL com a Emory University em uma temática interdisciplinar.
Marçal explicou que os professores da Univale têm doutorado em áreas tão distintas quanto Imunologia e Genética (Marçal); Odontologia (Elaine Pitanga) e Direito (Bernardo Barbosa).
Marília Ribeiro (Emory University) é doutora em Português e Estudos Brasileiros.
Os objetivos do projeto foram: desenvolver competências globais, como comunicação intercultural e trabalho em equipes internacionais; fortalecer a Internacionalização em Casa com atividades de ensino e pesquisa conjuntas; incentivar o desenvolvimento de projetos de pesquisa interdisciplinares (abordando temas como o papel da Literatura no Direito). continuação José Carlos Barbosa Lopes defendeu sua tese de doutorado em 2023 a partir de estudos sobre decolonização nos Intercâmbios Virtuais e de sua experiência prática no PCI que realiza desde o segundo semestre de 2020 entre Fatec Ipiranga e Jamestown Community College (JCC / SUNY, EUA).
Na seção “Artigo de Opinião”, neste número, leia as reflexões de Lopes sobre essa pesquisa.

% matched lemmas: 
\end{textsample}
